\subsection*{Implementation}

GraphPop is implemented as a Neo4j stored procedure plugin in Java (JDK~21, $\sim$6{,}000 lines). The 12 procedures are compiled into a single JAR file deployed to the Neo4j plugins directory. The import pipeline is a Python package (\texttt{graphpop-import}) using cyvcf2\cite{pedersen2017cyvcf2} for VCF parsing and numpy for vectorized genotype packing. VEP\cite{mclaren2016ensembl} consequence annotations and Reactome\cite{gillespie2022reactome} pathway data are imported via dedicated parsers. The MCP server (\texttt{graphpop-mcp}) is a Python package using the FastMCP framework, connecting to Neo4j via the official Python driver. All code is available under the MIT license.

\subsection*{Graph database}

GraphPop uses Neo4j Community Edition 2026.01.4. The graph schema comprises Variant, Sample, Population, Chromosome, Gene, Pathway, GOTerm, and GenomicWindow node labels with NEXT, HAS\_CONSEQUENCE, IN\_PATHWAY, HAS\_GO\_TERM, IN\_POPULATION, ON\_CHROMOSOME, and LD relationship types. Variant nodes are indexed on \texttt{(chr, pos)} for range queries. Population harmonic numbers ($a_n = H(2n-1)$, $a_{n2} = H_2(2n-1)$ for $n$ diploid samples) are stored on Population nodes and used by diversity and SFS procedures.

\subsection*{Packed genotype encoding}

Genotypes are stored as \texttt{gt\_packed} (2 bits/sample, encoding: 00=HomRef, 01=Het, 10=HomAlt, 11=Missing) and \texttt{phase\_packed} (1 bit/sample, which haplotype carries ALT). The cyvcf2 genotype encoding is remapped during import (0$\to$0, 1$\to$1, 2$\to$3, 3$\to$2) to match the internal 2-bit scheme. For non-diploid samples, \texttt{ploidy\_packed} (1 bit/sample) indicates ploidy; null implies all-diploid for backward compatibility. PackedGenotypeReader extracts genotype and phase values using bit-shift operations.

\subsection*{FAST PATH statistics}

\textbf{Diversity.} Nucleotide diversity $\pi = \sum_{i} 2p_i(1-p_i) \cdot n_i/(n_i-1)$ summed over variants, where $p_i = \text{AC}_i/\text{AN}_i$ and $n_i = \text{AN}_i/2$ for the focal population. Watterson's $\theta_W = S / a_n$ where $S$ is the number of segregating sites. Tajima's $D$ follows the standard formulation\cite{tajima1989statistical}. Fay \& Wu's $H = \pi - \theta_H$ where $\theta_H = \sum 2i^2/(n(n-1))$ is the squared-frequency estimator; the normalized version uses the Zeng et al.\ (2006) variance\cite{zeng2006statistical}.

\textbf{Divergence.} Hudson's Fst uses ratio-of-averages: $\text{Fst} = \sum N_i / \sum D_i$. Weir \& Cockerham (1984) Fst\cite{weir1984estimating} decomposes into variance components $(a, b, c)$; $\text{Fst} = \sum a / \sum(a+b+c)$. $D_{xy} = \sum p_1(1-p_2) + p_2(1-p_1)$. PBS (Population Branch Statistic)\cite{yi2010sequencing} uses $T = -\log(1 - \text{Fst}_\text{WC})$ for three population pairs.

\textbf{SFS.} Site frequency spectrum uses allele counts; unfolded SFS polarizes using \texttt{ancestral\_allele} annotations where available, with \texttt{n\_polarized} reporting coverage.

\subsection*{FULL PATH statistics}

\textbf{HaplotypeMatrix.} A dense bit-packed matrix loaded in one pass from Variant node arrays. Layout: \texttt{haplotypes[variantIdx][byteIdx]}, with haplotype $h$ accessed as \texttt{(haplotypes[v][h>>3] >> (h\&7)) \& 1}. Row length is $(\text{nHaplotypes}+7)/8$ bytes. \texttt{loadPair()} loads two populations from the same variant pass for XP-EHH.

\textbf{LD.} Pearson $r^2$ computed from haplotype dosages via VectorOps dot product. $D'$ computed from 2$\times$2 haplotype tables via \texttt{Integer.bitCount()} on packed byte arrays. LD edges are written between Variant nodes above a configurable $r^2$ threshold.

\textbf{EHH-based statistics.} EHHComputer walks position-sorted haplotype arrays from a focal variant outward, tracking distinct haplotype classes via flat arrays with incremental sum-of-pairs updates. iHH is the trapezoidal integral of the EHH decay curve. iHS\cite{voight2006map} $= \log(\text{iHH}_\text{ancestral}/\text{iHH}_\text{derived})$, standardized within allele-frequency bins (default 20). XP-EHH\cite{sabeti2007genome} $= \log(\text{iHH}_\text{pop1}/\text{iHH}_\text{pop2})$, genome-wide standardized. nSL\cite{ferrer2014defining} uses mean pairwise shared suffix length (SSL) per allele class. All EHH procedures are configurable: \texttt{min\_ehh} (truncation threshold), \texttt{max\_gap} (200~kb abort), \texttt{gap\_scale} (20~kb cap).

\textbf{Garud's $H$.}\cite{garud2015recent} Haplotype frequencies computed by hashing within sliding windows. $H_1 = \sum f_i^2$, $H_{12} = (f_1 + f_2)^2 + \sum_{i \geq 3} f_i^2$, $H_2/H_1 = (H_1 - f_1^2)/H_1$.

\textbf{ROH.} Two-state Viterbi HMM\cite{narasimhan2016estimating} with autozygous (AZ) and Hardy--Weinberg (HW) states. Emissions are AF-weighted: $P(\text{hom}|\text{AZ}) = 1 - \epsilon$, $P(\text{hom}|\text{HW}) = p^2 + (1-p)^2$ where $p$ is the population allele frequency. Transitions use continuous-time Markov chain (CTMC) matrix power for inter-variant distances, matching bcftools roh\cite{narasimhan2016bcftools}. Default parameters: \texttt{hw\_to\_az}~$= 6.7 \times 10^{-8}$, \texttt{az\_to\_hw}~$= 5 \times 10^{-9}$, \texttt{het\_error\_rate}~$= 10^{-3}$.

\subsection*{VariantFilter and VariantQuery}

All procedures share a common filtering infrastructure. VariantFilter supports \texttt{min\_af}, \texttt{max\_af}, \texttt{min\_call\_rate}, \texttt{hwe\_pvalue} (Wigginton et al.\ 2005 exact test\cite{wigginton2005note}), and \texttt{variant\_type}. VariantQuery dynamically constructs Cypher queries that optionally join through \texttt{HAS\_CONSEQUENCE} edges (for consequence filtering) or \texttt{IN\_PATHWAY} edges (for pathway filtering) before the main variant scan. Custom sample subsets are supported via \texttt{samples}/\texttt{samples1}/\texttt{samples2} options using SampleSubsetComputer for FAST PATH or HaplotypeMatrix.load() overloads for FULL PATH.

\subsection*{Datasets}

\textbf{1000 Genomes Project.} Phase~3 chr22 VCF (3{,}202 samples, 1{,}070{,}401 variants, 5 superpopulations: AFR, AMR, EAS, EUR, SAS)\cite{1000genomes2015global}. Used for benchmarking and validation. Full-genome analysis across all autosomes was also performed.

\textbf{3K Rice Genomes Project.} 3{,}024 accessions, $\sim$29.6M SNPs, 12 chromosomes, 12 subpopulations (Wang et al.\ 2018 ADMIXTURE K=9 classification)\cite{wang2018genomic}. VEP annotations with the \emph{Oryza sativa} IRGSP-1.0 reference. Ancestral allele annotations derived from EPO multi-species alignment FASTA.

\subsection*{Benchmark methodology}

Each benchmark was run 3 times with 1 warm-up run discarded. GraphPop was called via the Neo4j Python driver with the database pre-loaded in memory (warm cache). Competitor tools read from their native file formats (VCF for VCFtools/scikit-allel/bcftools, PLINK binary for PLINK~2/1.9). Peak memory was measured using GraphPop's \texttt{graphpop.heap\_stats()} procedure (JVM heap) and \texttt{/usr/bin/time -v} (RSS) for external tools. Runtime comparisons used wall-clock time including I/O.

\subsection*{Hardware}

All benchmarks and analyses were performed on a single workstation running WSL2 on Windows, with an AMD processor, 64~GB RAM, and NVMe SSD storage. Neo4j was configured with 4~GB heap and 16~GB page cache.
